\section{USAMO 1992}
        \begin{enumerate}
        	\item (\textit{USAMO 1992}) Find, as a function of $\, n, \,$ the sum of the digits of 

 \[9 \times 99 \times 9999 \times \cdots \times \left( 10^{2^n} - 1 \right),\] 
where each factor has twice as many digits as the previous one.


 



\textbf{Solution}

The answer is $9 \cdot 2^n$.

 Let us denote the quantity $\prod_{k=0}^n \bigl( 10^{2^k}-1 \bigr)$ as $P_n$.  We wish to find the sum of the digits of $P_n$.

 We first note that
 \[P_{n-1} < \prod_{k=0}^{n-1} 10^{2^k} = 10^{2^n-1},\]so $P_{n-1}$ is a number of at most $2^n$ digits.  We also note that the units digit is not equal to zero.  We may thus represent $P_{n-1}$ as
 \[\sum_{k=0}^{2^n-1} 10^k d_k ,\]where the $d_k$ are digits and $d_0 \neq 0$.  Then
 \begin{align*} P_n &= \bigl( 10^{2^n}-1 \bigr) P_{n-1} = \sum_{k=0}^{2^n-1} - 10^k d_k + \sum_{k=0}^{2^n-1} 10^{2^n+k} d_k \\ &= (10-d_0) + \sum_{k=1}^{2^n-1} 10^k(9-d_k) + 10^{2^n}(d_0-1) + \sum_{k=1}^{2^n-1} 10^{2^n+k} d_k . \end{align*}
Thus the digits of $P_n$ are
 \[10-d_0, 9-d_1, 9-d_2, \dotsc, 9-d_{2^n-1}, d_0-1, d_1, d_2, \dotsc, d_{2^n-1} ,\]and the sum of these is evidently $9 \cdot 2^n$, as desired.  $\blacksquare$

 
<i>Alternate solutions are always welcome. If you have a different, elegant solution to this problem, please \href{https://artofproblemsolving.com/wiki/index.php?title=1992\_USAMO\_Problems/Problem\_1&action=edit}{add it} to this page.</i>

 


	\item (\textit{USAMO 1992}) Prove 

 \[\frac{1}{\cos 0^\circ \cos 1^\circ} +  \frac{1}{\cos 1^\circ \cos 2^\circ} + \cdots + \frac{1}{\cos 88^\circ \cos 89^\circ} = \frac{\cos 1^\circ}{\sin^2 1^\circ}.\]


 



\textbf{Solution 1}

Consider the points $M_k = (1, \tan k^\circ)$ in the coordinate plane with origin $O=(0,0)$, for integers $0 \le k \le 89$.

 
                    \begin{figure}[H]
                        \centering
                        \includegraphics[width=0.5\textwidth]{Solution/USAMO_I_1992_2_Solution_Figure_1.png}
                    \end{figure}
                    

 

 Evidently, the angle between segments $OM_a$ and $OM_b$ is $(b-a)^\circ$, and the length of segment $OM_a$ is $1/\cos a^\circ$.  It then follows that the area of triangle $M_aOM_b$ is $\frac{1}{2} \sin(b-a)^\circ \cdot OM_a \cdot OM_b = \frac{1}{2} \sin(b-a)^\circ / (\cos a^\circ \cdot \cos b^\circ)$.  Therefore
 \begin{align*} \sum_{k=0}^{88} \frac{\frac{1}{2} \sin 1^\circ}{ \cos k^\circ \cos k+1^\circ} &= \sum_{k=0}^{88} [ M_k O M_{k+1} ] \\ &= [ M_0 O M_{89} ] \\ &= \frac{ \frac{1}{2} \sin 89^\circ}{\cos 89^\circ} = \frac{\frac{1}{2} \cos 1^\circ}{\sin 1^\circ} , \end{align*}
so
 \[\sum_{k=0}^{88} \frac{1}{\cos k^\circ \cos (k+1)^\circ} = \frac{\cos 1^\circ}{ \sin^2 1^\circ},\]as desired.  $\blacksquare$

 



\textbf{Solution 2}

First multiply both sides of the equation by $\sin 1$, so the right hand side is $\frac{\cos 1}{\sin 1}$.  Now by rewriting $\sin 1=\sin((k+1)-k)=\sin(k+1)\cos(k)+\sin(k)\cos(k+1)$, we can derive the identity $\tan(n+1)-\tan(n)=\frac{\sin 1}{\cos(n)\cos(n+1)}$.  Then the left hand side of the equation simplifies to $\tan 89-\tan 0=\tan 89=\frac{\sin 89}{\cos 89}=\frac{\cos 1}{\sin 1}$ as desired.

 



\textbf{Solution 3}

Multiply by $\sin{1}$. We get:

 $\frac {\sin{1}}{\cos{0}\cos{1}} + \frac {\sin{1}}{\cos{1}\cos{2}} + ... + \frac {\sin{1}}{\cos{88}\cos{89}} = \frac {\cos{1}}{\sin{1}}$

 we can write this as:

 $\frac {\sin{1 - 0}}{\cos{0}\cos{1}} + \frac {\sin{2 - 1}}{\cos{1}\cos{2}} + ... + \frac {\sin{89 - 88}}{\cos{88}\cos{89}} = \frac {\cos{1}}{\sin{1}}$

 This is an identity $\tan{a} - \tan{b} = \frac {\sin{(a - b)}}{\cos{a}\cos{b}}$

 Therefore;

 $\sum_{i = 1}^{89}[\tan{k} - \tan{(k - 1)}] = \tan{89} - \tan{0} = \cot{1}$, because of telescoping. 

 but since we multiplied $\sin{1}$ in the beginning, we need to divide by $\sin{1}$. So we get that:

 $\frac {1}{\cos{0}\cos{1}} + \frac {1}{\cos{1}\cos{2}} + \frac {1}{\cos{2}\cos{3}} + .... + \frac {1}{\cos{88}\cos{89}} = \frac {\cos{1}}{\sin^2{1}}$as desired. QED

 



\textbf{Solution 4}

Let $S = \frac{1}{\cos 0^\circ\cos 1^\circ} + \frac{1}{\cos 1^\circ\cos 2^\circ} + ... + \frac{1}{\cos 88^\circ\cos 89^\circ}$.

 Multiplying by $\sin 1^\circ$ gives
 \[S \sin 1^\circ = \frac{\sin(1^\circ-0^\circ)}{\cos 0^\circ\cos 1^\circ} + ... + \frac{\sin(89^\circ-88^\circ)}{\cos 88^\circ\cos 89^\circ}\]

 Notice that $\frac{\sin((x+1^\circ)-x)}{\cos 0^\circ\cos 1^\circ} = \tan (x+1^\circ) - \tan x$ after expanding the sine, and so
 \[S \sin 1^\circ = \left(\tan 1^\circ - \tan 0^\circ\right) + \cdots + \left(\tan 89^\circ - \tan 88^\circ\right) = \tan 89^\circ - \tan 0^\circ = \cot 1^\circ = \frac{\cos 1^\circ}{\sin 1^\circ},\] so 
 \[S = \frac{\cos 1^\circ}{\sin^21^\circ}.\]

 


	\item (\textit{USAMO 1992}) For a nonempty set $S$ of integers, let $\sigma(S)$ be the sum of the elements of $S$. Suppose that $A = \{a_1, a_2, \ldots, a_{11}\}$ is a set of positive integers with $a_1 < a_2 < \cdots < a_{11}$ and that, for each positive integer $n \le 1500$, there is a subset $S$ of $A$ for which $\sigma(S) = n$. What is the smallest possible value of $a_{10}$?


 



\textbf{Solution}

Let's a $n$-$p$ set be a set $Z$ such that $Z=\{a_1,a_2,\cdots,a_n\}$, where $\forall i<n$, $i\in \mathbb{Z}^+$, $a_i<a_{i+1}$, and for each $x\le p$, $x\in \mathbb{Z}^+$, $\exists Y\subseteq Z$, $\sigma(Y)=x$, $\nexists \sigma(Y)=p+1$.

 (For Example $\{1,2\}$ is a $2$-$3$ set and $\{1,2,4,10\}$ is a $4$-$8$ set)

 
Furthermore, let call a $n$-$p$ set a $n$-$p$ good set if $a_n\le p$, and a $n$-$p$ bad set if $a_n\ge p+1$ (note that $\nexists \sigma(Y)=p+1$ for any $n$-$p$ set. Thus, we can ignore the case where $a_n=p+1$).

 
Furthermore, if you add any amount of elements to the end of a $n$-$p$ bad set to form another $n$-$p$ set (with a different $n$), it will stay as a $n$-$p$ bad set because $a_{n+x}>a_{n}>p+1$ for any positive integer $x$ and $\nexists \sigma(Y)=p+1$.

 
\textbf{Lemma}) If $Z$ is a $n$-$p$ set, $p\leq 2^n-1$.

 
For $n=1$, $p=0$ or $1$ because $a_1=1 \rightarrow p=1$ and $a_1\ne1\rightarrow p=0$.

 Assume that the lemma is true for some $n$, then $2^n$ is not expressible with the $n$-$p$ set. Thus, when we add an element to the end to from a $n+1$-$r$ set, $a_{n+1}$ must be $\le p+1$ if we want $r>p$ because we need a way to express $p+1$. Since $p+1$ is not expressible by the first $n$ elements, $p+1+a_{n+1}$ is not expressible by these $n+1$ elements. Thus, the new set is a $n+1$-$r$ set, where $r\leq p+1+a_{n+1} \leq 2^{n+1}-1$

 \textbf{Lemma Proven}

 
The answer to this question is $\max{(a_{10})}=248$.

 The following set is a $11$-$1500$ set: 

 $\{1,2,4,8,16,32,64,128,247,248,750\}$

 Note that the first 8 numbers are power of $2$ from $0$ to $7$, and realize that any $8$ or less digit binary number is basically sum of a combination of the first $8$ elements in the set. Thus, $\exists Y\subseteq\{1,2,4,8,16,32,64,128\}$, $\sigma(Y)=x \forall 1\le x\leq 255$. 

 $248\le\sigma(y)+a_9\le502$ which implies that $\exists A\subseteq\{1,2,4,8,16,32,64,128,247\}$, $\sigma(A)=x \forall 1\le x\leq 502$.

 Similarly $\exists B\subseteq\{1,2,4,8,16,32,64,128,247,248\}$, $\sigma(A)=x \forall 1\le x\le750$ and $\exists C\subseteq\{1,2,4,8,16,32,64,128,247,248,750\}$, $\sigma(A)=x \forall 1\leq x\leq 1500$.

 
Thus, $\{1,2,4,8,16,32,64,128,247,248,750\}$ is a $11$-$1500$ set.

 

Now, let's assume for contradiction that $\exists a_{10}\leq 247$ such that ${a_1, a_2, \dots, a_{11}}$ is a $11$-$q$ set where $q\geq 1500$

 ${a_1, a_2, \dots a_8}$ is a $8$-$a$ set where $a\leq 255$ (lemma).

 $max{(a_9)}=a_{10}-1\leq 246$

 Let ${a_1, a_2, \dots, a_{10}}$ be a $10$-$b$ set where the first $8$ elements are the same as the previous set. Then, $256+a_9+a_{10}$ is not expressible as $\sigma(Y)$. Thus, $b\leq 255+a_9+a_{10}\leq 748$.

 In order to create a $11$-$d$ set with $d>748$ and the first $10$ elements being the ones on the previous set, $a_{11}\leq 749$ because we need to make $749$ expressible as $\sigma(Y)$. Note that $b+1+a_{11}$ is not expressible, thus $d<b+1+a_{11}\leq 1498$.

 
Done but not elegant...

 


	\item (\textit{USAMO 1992}) Chords $AA'$, $BB'$, and $CC'$ of a sphere meet at an interior point $P$ but are not contained in the same plane.  The sphere through $A$, $B$, $C$, and $P$ is tangent to the sphere through $A'$, $B'$, $C'$, and $P$.  Prove that $AA'=BB'=CC'$.


 



\textbf{Solution 1}

Consider the plane through $A,A',B,B'$.  This plane, of course, also contains $P$. We can easily find the $\triangle APB$ is isosceles because the base angles are equal. Thus, $AP=BP$. Similarly, $A'P=B'P$. Thus, $AA'=BB'$. By symmetry, $BB'=CC'$ and $CC'=AA'$, and hence $AA'=BB'=CC'$ as desired. 

 $\mathbb{QED.}$

 \textbf{Add-on} By another person \textasciicircum{}v\textasciicircum{}

 The person that came up with the solution did not prove that  $\triangle APB$ is isosceles nor the base angles are congruent. I will add on to the solution.

 There is a common tangent plane that pass through $P$ for the $2$ spheres that are tangent to each other.

 
Since any cross section of sphere is a circle. It implies that $A$, $A'$, $B$, $B'$ be on the same circle ($\omega_1$), $A$, $B$, $P$ be on the same circle ($\omega_2$), and $A'$, $B'$, $P$ be on the same circle ($\omega_3$).

 $m\angle APB= m\angle A'PB'$ because they are vertical angles. By power of point, $(AP)(A'P)=(BP)(B'P)\rightarrow\frac{AP}{B'P}=\frac{BP}{A'P}$

 By the SAS triangle simlarity theory, $\triangle APB \sim\triangle B'PA'$. That implies that $\angle ABP\cong\angle B'PA'$.

 


 Let's call the interception of the common tangent plane and the plane containing $A$, $A'$, $B$, $B'$, $P$, line $l$.

 $l$ must be the common tangent of $\omega_2$ and $\omega_3$.

 The acute angles form by $l$ and $\overline{AA'}$ are congruent to each other (vertical angles) and by the tangent-chord theorem, the central angle of chord $\overline{AP}$ and $\overline{A'P}$ are equal.

 Similarly the central angle of chord $\overline{BP}$ and $\overline{B'P}$ are equal.

 The length of any chord with central angle $2\theta$ and radius $r$ is $2r\sin\left({\theta}\right)$, which can easily been seen if we drop the perpendicular from the center to the chord.

 Thus, $\frac{AP}{A'P}=\frac{BP}{B'P}$.

 By the SAS triangle simlarity theory, $\triangle APB \sim\triangle A'PB'$. That implies that $\angle ABP\cong\angle B'PA'$.

 
That implies that $\angle ABP\cong\angle A'PB'\cong\angle B'PA'$. Thus, $\triangle A'PB'$ is an isosceles triangle and since $\triangle APB \sim\triangle A'PB'$, then$\triangle APB$ must be an isosceles triangle too.

 



\textbf{Solution 2}

Call the large sphere $O_1$, the one containing $A$ $O_2$, and the one containing $A'$ $O_3$.  The centers are $c_1$, $c_2,$ and $c_3$.

 
Since two spheres always intersect in a circle , we know that A,B, and C must lie on a circle ($w_1$)completely contained in $O_1$ and $O_2$

 
Similarly, A', B', and C' must lie on a circle ($w_2$) completely contained in $O_1$ and $O_3$.  

 
So, we know that 3 lines connecting a point on $w_1$ and P hit a point on $w_2$.  This implies that $w_1$ projects through P to $w_2$, which in turn means that $w_1$ is in a plane parallel to that of $w_2$.  Then, since $w_1$ and $w_2$ lie on the same sphere, we know that they must have the same central axis, which also must contain P (since the center projects through P to the other center).  

 
So, all line from a point on $w_1$ to P are of the same length, as are all lines from a point on $w_2$ to P.  Since AA', BB', and CC' are all composed of one of each type of line, they must all be equal.

 


	\item (\textit{USAMO 1992}) Let $P(z)$be a polynomial with complex coefficients which is of degree $1992$ and has distinct zeros.Prove that there exists complex numbers $a_1,a_2,\cdots,a_{1992}$ such that $P(z)$ divides the polynomial


 $\left(\cdots\left(\left(z-a_1\right)^2-a_2\right)^2\cdots-a_{1991}\right)^2-a_{1992}$

 



\textbf{Solution}

Since the zeros $z_{1}, z_{2}, \ldots, z_{1992}$ of the polynomial $P(z)$ are distinct, the polynomial $P(z)$ divides the polynomial 
 \[Q\left(z\right) = \left( \ldots \left( \left(z-a_{1}\right)^{2}-a_{2}\right)^{2}\ldots-a_{1991}\right)^{2}-a_{1992}\]if and only if $Q\left(z_{i}\right)=0$ for every $i\in\left\{1;\;2;\;\ldots;\;1992\right\}$. So it is enough to show that for any $1992$ complex numbers $z_{1}, z_{2}, \dots, z_{1992}$, there exist $1992$ complex numbers $a_{1}, a_{2}, \ldots, a_{1992}$, such that the polynomial 
 \[Q\left(z\right) = \left(\ldots\left( \left(z-a_{1}\right)^{2}-a_{2}\right)^{2}\ldots-a_{1991}\right)^{2}-a_{1992}\]satisfies $Q\left(z_{i}\right)=0$ for every $i\in\left\{1;\;2;\;\ldots;\;1992\right\}$. This is a special case of the following lemma.

 
\textbf{Lemma:} Let $n$ be a positive integer, and $z_{1}, z_{2}, \ldots, z_{n}$ be $n$ complex numbers. Then, there exist $n$ complex numbers $a_{1}, a_{2},\ldots, a_{n}$ such that the polynomial $Q\left(z\right) = \left( \ldots \left( \left(z-a_{1}\right)^{2}-a_{2}\right)^{2}\ldots-a_{n-1}\right)^{2}-a_{n}$ satisfies $Q\left(z_{i}\right)=0$ for every $i\in\left\{1;\;2;\;\ldots;\;n\right\}$.

 
\textbf{Proof:} We use induction over $n$.

 For $n = 1$, the lemma is trivial, since $Q\left(z\right)=z-a_{1}$, so we can take $a_{1}=z_{1}$, and then the polynomial $Q\left(z\right)=z-z_{1}$ clearly satisfies $Q\left(z_{1}\right)=0$.

 
Let $k$ be a positive integer. Assume that for $n = k$, the lemma is true. Therefore, there exist $k$ complex numbers $a_{1}, a_{2}, \ldots, a_{k}$ such that the polynomial $Q\left(z\right)=\left(\ldots\left( \left(z-a_{1}\right)^{2}-a_{2}\right)^{2}\ldots-a_{k-1}\right)^{2}-a_{k}$ satisfies $Q\left(z_{i}\right)=0$ for every $i\in\left\{1;\;2;\;\ldots;\;k\right\}$. Then we wish to prove that for any $k + 1$ complex numbers $z_{1}, z_{2}, \ldots, z_{k+1}$, there exist $k + 1$ complex numbers $b_{1}, b_{2}, \ldots, b_{k+1}$ such that the polynomial $R\left(z\right)=\left(\ldots\left( \left(z-b_{1}\right)^{2}-b_{2}\right)^{2}\ldots-b_{k}\right)^{2}-b_{k+1}$ satisfies $R\left(z_{i}\right)=0$ for every $i\in\left\{1;\;2;\;\ldots;\;k+1\right\}$. We will construct this polynomial $R$ from the polynomial $Q$.

 Let $b_{i}=a_{i}$ for every $i\le k-1$, and let $b_{k}=a_{k}+\frac{1}{2}Q\left(z_{k+1}\right)$ and $b_{k+1}=\frac{1}{4}\left(Q\left(z_{k+1}\right)\right)^{2}$. Then,

 

 \begin{align*} R\left(z\right)&=\left(\ldots\left(\left(\left(z-b_{1}\right)^{2}-b_{2}\right)^{2}\ldots-b_{k-1}\right)^{2}-b_{k}\right)^{2}-b_{k+1}\\ & =\left(\ldots\left(\left(\left(z-a_{1}\right)^{2}-a_{2}\right)^{2}\ldots-a_{k-1}\right)^{2}-\left(a_{k}+\frac{1}2Q\left(z_{k+1}\right)\right)\right)^{2}-\frac{1}4\left(Q\left(z_{k+1}\right)\right)^{2} \\ & =\left(\left(\ldots\left(\left(\left(z-a_{1}\right)^{2}-a_{2}\right)^{2}\ldots-a_{k-1}\right)^{2}-a_{k}\right)-\frac{1}2Q\left(z_{k+1}\right)\right)^{2}-\frac{1}4\left(Q\left(z_{k+1}\right)\right)^{2}\\ &=\left(Q\left(z\right)-\frac{1}2Q\left(z_{k+1}\right)\right)^{2}-\frac{1}4\left(Q\left(z_{k+1}\right)\right)^{2}=Q\left(z\right)\left(Q\left(z\right)-Q\left(z_{k+1}\right)\right). \end{align*}


 
Then as $Q\left(z_{i}\right)=0$ for $i\in\left\{1;\;2;\;\ldots;\;k\right\}$, we have 
 \[R\left(z_{i}\right)=Q\left(z_{i}\right)\left(Q\left(z_{i}\right)-Q\left(z_{k+1}\right)\right)=0\]for all $i\in\left\{1;\;2;\;\ldots;\;k\right\}$. Also, 
 \[R\left(z_{k+1}\right)=Q\left(z_{k+1}\right)\left(Q\left(z_{k+1}\right)-Q\left(z_{k+1}\right)\right)=0.\]Thus, $R\left(z_{i}\right)=0$ holds for every $i\in\left\{1;\;2;\;\ldots;\;k+1\right\}$, and this proves the lemma for $n = k + 1$. Hence, the induction step is complete, and the lemma is proven.

 As the problem is a special case of the lemma ($n=1992$), we are done.

 


\end{enumerate}